\begin{lemma}[Four common missing colors at an uncolored edge]
Let $G=(V,E)$ be a finite loopless multigraph, let $e_0\in E$, and let
$\varphi$ be a partial $k$-edge-coloring which is proper and leaves exactly
$e_0$ uncolored.  Suppose
\[
                    \Delta(G)+1\leq k.
\]
Let $A$ and $B$ be tree-sequences with distinguished first edge $e_0$.
Write $\partial e_0=\{x,y\}$, and suppose that both $x$ and $y$ belong to
both $V(A)$ and $V(B)$.  If $V(A)\cup V(B)$ is elementary under
$\varphi$, then
\[
 \bigl|\overline\varphi(A)\cap\overline\varphi(B)\bigr|\geq 4.
\]
Here tree-missing colors are understood in the sense of
\noderef{TreeMissingColors}, so the intersection is an intersection of
finite sets of exact color identities.
\end{lemma}
\begin{proof}
Because the loopless law in \noderef{LooplessMultigraph} excludes an edge
with two equal ends, and the two ends of $e_0$ are $x$ and $y$, one has
$x\ne y$.  We first prove that at least two colors are missing at each end
of $e_0$.

Fix an end $v\in\{x,y\}$, and put $d=d_G(v)$.  Consider the finite set
$I_v$ of exact edge identities incident with $v$ which are colored by
$\varphi$.  By the hypothesis that exactly $e_0$ is uncolored, as defined
in \noderef{IsOneEdgeUncolored},
\[
                    I_v=\{e:e\text{ is incident with }v,\ e\ne e_0\}.
\]
The endpoint equation and \noderef{Incident} say that $e_0$ is incident
with $v$, so the definition of degree in \noderef{Degree} gives
\[
                         |I_v|=d-1.                              \tag{1}
\]

Sending an edge of $I_v$ to its assigned color is injective.  Indeed, if
two distinct identities in $I_v$ received the same color, they would be
adjacent because they share the vertex $v$, by
\noderef{AdjacentEdges}; this would contradict properness as defined in
\noderef{IsProperEdgeColoring}.  Conversely, the witness clause in the
definition of \noderef{PresentColors} says that every color present at
$v$ is the color of some identity in $I_v$.  Thus the injection is a
bijection from $I_v$ onto the set of colors present at $v$, and (1) yields
\[
             |\varphi(v)|=d-1,
             \qquad |\overline\varphi(v)|=k-(d-1).              \tag{2}
\]
The second equality is precisely the complement definition in
\noderef{MissingColors}.
The definition of \noderef{MaximumDegree} gives $d\leq\Delta(G)$.
Together with $\Delta(G)+1\leq k$, equation (2) implies
\[
                         |\overline\varphi(v)|\geq 2.            \tag{3}
\]
Applying (3) to $x$ and to $y$ gives two missing-color sets, each of
cardinality at least two.

Both $x$ and $y$ lie in $V(A)\cup V(B)$, and they are distinct.  The
elementarity hypothesis, in the precise sense of
\noderef{ColoredVertexSetElementary}, therefore gives
\[
             \overline\varphi(x)\cap\overline\varphi(y)
             =\varnothing.                                     \tag{4}
\]
Since each of $x,y$ belongs to $V(A)$, the union definition in
\noderef{TreeMissingColors} gives
\(
 \overline\varphi(x)\cup\overline\varphi(y)
 \subseteq\overline\varphi(A)
\).
The same argument using membership in $V(B)$ gives containment in
$\overline\varphi(B)$.  Hence
\[
 \overline\varphi(x)\cup\overline\varphi(y)
 \subseteq
 \overline\varphi(A)\cap\overline\varphi(B).                   \tag{5}
\]
By (3) and the disjointness (4), the set on the left of (5) has at least
$2+2=4$ elements.  Taking finite-set cardinalities in (5) proves the
claim.
\end{proof}
